\documentclass[11pt]{ctexart}
\usepackage[a4paper,margin=2.2cm]{geometry}
\usepackage{amsmath,amssymb}
\usepackage{booktabs}
\usepackage{fancyhdr}
\pagestyle{fancy}
\fancyhf{}
\fancyhead[L]{2026 Fall: High-dimensional Data Analysis}
\fancyhead[R]{自测版（无答案）}
\fancyfoot[C]{\thepage}
\renewcommand{\headrulewidth}{0.4pt}

\newcommand{\fillblank}{\mbox{（\qquad）}}

\title{\textbf{Sample Midterm Exam 1}\\[0.2em]
\large （自测版 · 无答案）}
\author{}
\date{}

\begin{document}
\maketitle
\vspace{-1.2cm}

\noindent\textbf{说明：}本样卷供同学们参考。正式考试中的风格和难度与此样卷类似。

\vspace{0.6em}

\noindent\textbf{Problem 1 (10 pts).} 判断题，请判断以下说法为正确 (T) 或者错误 (F)。
\begin{enumerate}
\item 简单线性回归中，若用最小二乘 (OLS) 估计，则数据的中心点 $(\bar{x}, \bar{y})$ 一定落在估计出的回归线上。\fillblank
\item 在包含 $p$ 个自变量的线性回归 $Y \sim X_1 + \cdots + X_p$ 中，用最小二乘法估计回归系数，则 $\hat{\sigma}^2 = \mathrm{RSS}/(n-p-1)$ 是 $\sigma_*^2$ 的无偏估计。\fillblank
\item 对于一个有 $d$ 个水平的类别型变量 $F$，做线性回归 $Y \sim F$。若采用虚拟变量法引入 $(d-1)$ 个虚拟变量，则所有虚拟变量的回归系数都代表该类别与 $Y$ 的总样本均值 $\bar{y} = \sum_{i=1}^n y_i / n$ 的差异。\fillblank
\item 在 logistic 回归，回归的 $G^2$ (deviance) 越小，说明模型拟合得越好。\fillblank
\item 非线性回归的 Gauss-Newton 算法一定保证收敛到全局最优解。\fillblank
\end{enumerate}

\noindent\textbf{Problem 2 (16 pts).} 单选题，每题只有一个正确答案。
\begin{enumerate}
\item 在多元线性回归中，加入一个与因变量完全不相关但与其他自变量高度相关的变量，最可能导致\fillblank\\
(A) 新变量的系数会非常显著（即 $p$-值很小）\quad
(B) 所有变量系数的标准误都会减小\\
(C) 原来自变量的系数出现变化，甚至由显著变得不显著\quad
(D) 模型的 $R^2$ 显著提高
\item 以下关于 $R^2_{\mathrm{adj}}$ 的说法中，正确的是\fillblank\\
(A) 其唯一目标是提升模型在训练数据上的拟合优度\quad
(B) 它可以帮助我们避免过拟合\\
(C) 它确保了模型中所有变量的系数都是统计显著的\quad
(D) 当选取的变量数 $k$ 固定时，$R^2_{\mathrm{adj}}$ 就完全确定了
\item 关于 lasso 和岭回归的说法，错误的是\fillblank\\
(A) lasso 可以将某些系数压缩至 $0$，从而实现变量选择；而岭回归通常做不到\\
(B) 两者都可以用于处理变量数量多于样本数量（$p > n$）的情况\\
(C) 两者估计出的回归系数都是有偏的\\
(D) 对于每一个自变量，lasso 生成的回归系数的长度（即绝对值）一定比岭回归要小
\item 关于主成分分析，以下说法错误的是\fillblank\\
(A) 主成分是原始变量的线性组合，且各主成分之间相互正交\\
(B) 通常需要保留所有特征值大于 $0$ 的主成分\\
(C) 载荷表示主成分中每个变量的权重\\
(D) 得分表示每个样本点在主成分方向上的值
\end{enumerate}

\noindent\textbf{Problem 3 (15 pts).} 在一个简单线性回归中，样本数量为 $n = 14$，计算出 $\mathrm{SYY} = 20$，$\mathrm{SXX} = 20$，$\mathrm{SXY} = 10$，$\mathrm{RSS} = 15$。
\begin{enumerate}
\item (5 pts) 斜率的最小二乘估计 $\hat{\beta}_1$。
\item (5 pts) 求出可决系数 $R^2$。
\item (5 pts) 构造 $t$-检验来判断 $\beta_1^*$ 是否显著非 $0$。这里显著水平选 $\alpha = 0.05$，已知 $t(12)$-分布的上分位数 $t_{\alpha/2}(12) = 2.45$。
\end{enumerate}

\noindent\textbf{Problem 4 (14 pts).} 假设一个线性回归中，有一个连续自变量为 $X$ 和一个分类自变量 $F$，$F$ 有 $3$ 个类别。样本量为 $n = 10$。我们从一个只有截距的模型开始，用向前选变量法做了三步，其中 $F$ 生成的虚拟变量整组一起加进来。表~\ref{tab:mt1p4} 给出了过程中拟合的四个模型以及相应的 RSS。注意模型 (2), (3), (4) 也都自然包含了截距。

\begin{table}[h]
\centering
\caption{Problem 4 的四个模型}
\label{tab:mt1p4}
\begin{tabular}{clc}
\toprule
序号 & 模型 & RSS \\
\midrule
(1) & $Y \sim 1$（只有截距的模型） & 100 \\
(2) & $Y \sim X$ & 48 \\
(3) & $Y \sim X + F$ & 24 \\
(4) & $Y \sim X + F + X : F$ & 18 \\
\bottomrule
\end{tabular}
\end{table}

\begin{enumerate}
\item (7 pts) 我们知道，adjusted $R^2$ 的定义为
\[
R^2_{\mathrm{adj}} = 1 - \frac{n-1}{n-k-1}\left(1 - R^2\right),
\]
其中 $k$ 为模型中包含的自变量个数。试用此准则选出合适的模型。
\item (7 pts) 已知 AIC 的定义为
\[
\mathrm{AIC} = n \log\left(\frac{\mathrm{RSS}}{n}\right) + 2(k+1).
\]
试用此准则选出合适的模型。（已知：$e^{0.2} \approx 1.2$，$e^{0.4} \approx 1.5$。）
\end{enumerate}

\noindent\textbf{Problem 5 (15 pts).} 一次商业推广活动中，商家统计了一些用户的基本信息，以及每位用户是否最后有购买行为。我们通过 logistic 回归来分析用户的购买行为受哪些因素影响。我们选用的数据样本量为 $n = 500$，每位用户的信息包括
\begin{itemize}
\item $X$：收入（万元），为连续变量；
\item $A$：性别，为分类变量，分为“女性 (F)”、“男性 (M)”；
\item $Y$：是否购买，为 $0/1$ 变量。若存在购买行为，则取值为 $1$；反之取值为 $0$。
\end{itemize}
现在拟合一个 logistic 回归模型 $Y \sim X + A + X : A$，得到表~\ref{tab:mt1p5} 中的 summary 结果。

\begin{table}[h]
\centering
\caption{Summary table}
\label{tab:mt1p5}
\begin{tabular}{lrrrr}
\toprule
 & Estimate & Std.\ Error & z value & $\Pr(>|z|)$ \\
\midrule
(Intercept) & 1.55 & 1.20 & 1.32 & 0.19 \\
$X$ & 0.47 & 0.22 & 2.11 & 0.035 \\
$M$ & $-1.66$ & 0.53 & $-3.15$ & 0.0016 \\
$X : M$ & 0.35 & 0.40 & 0.88 & 0.38 \\
\bottomrule
\end{tabular}
\end{table}

\noindent Null Deviance: 188.50 on 499 degrees of freedom\\
Residual deviance: 43.12 on 496 degrees of freedom

\begin{enumerate}
\item (5 pts) 对于一个收入为 $x_{\mathrm{new}} = 1.5$ 的女性新用户，试用表中的数据给出其购买的概率。可以只列式不计算。
\item (5 pts) 我们比较女性用户与男性用户的购买 odds（即购买概率/不购买概率）。固定收入水平 $x = 1$，构建一个全模型和一个主效应模型。在这两个模型之下，分别说明两性用户的购买 odds 之间的关系。
\item (4 pts) 在未来的商业推广活动中，对用户群体购买行为的预测决定了营销策略的方向。经理 A 说：“我们应该不要漏掉可能的客户，即使这样会使我们在未来的商业推广活动中产生更高的费用。”经理 B 说：“我们应该更专注特定的用户群体进行精准营销，这样可以节省营销成本。”针对两个经理的话，我们分别应该怎么去预测下次营销活动中新用户的购买行为？
\end{enumerate}

\noindent\textbf{Problem 6 (16 pts).} 对于线性回归，我们令 $A = (X^{\top}X)^{-1}X^{\top} \in \mathbb{R}^{(p+1)\times n}$，则最小二乘估计量可以写为 $\hat{\beta} = AY$。现在给定任何一个满足 $\tilde{A}X = I$ 的常数矩阵 $\tilde{A} \in \mathbb{R}^{(p+1)\times n}$，考虑一个估计量 $\tilde{\beta} = \tilde{A}Y$。
\begin{enumerate}
\item (4 pts) 证明：$\tilde{\beta}$ 为 $\beta^*$ 的无偏估计量。
\item (4 pts) 对于任何常数向量 $c \in \mathbb{R}^{p+1}$，令 $u_1 = c^{\top}\hat{\beta}$，$u_2 = c^{\top}\tilde{\beta}$。证明：$u_1, u_2$ 均为 $c^{\top}\beta^*$ 的无偏估计量。
\item (4 pts) 证明：$\mathrm{Cov}(u_1, u_2) = \mathrm{Var}(u_1)$。
\item (4 pts) 利用 3 的结论，证明 $\mathrm{Var}(u_1) \leq \mathrm{Var}(u_2)$。
\end{enumerate}

\noindent\textbf{Problem 7 (15 pts).} 假设两个自变量 $Z_1, Z_2$ 分别有一列样本值 $Z_1 = (z_{11}, \ldots, z_{n1})^{\top} \in \mathbb{R}^n$，$Z_2 = (z_{12}, \ldots, z_{n2})^{\top} \in \mathbb{R}^n$，它们都已中心化、标准化，并且互相正交。即 $\sum_{i=1}^n z_{i1}/n = \sum_{i=1}^n z_{i2}/n = 0$，$\|Z_1\|^2 = \|Z_2\|^2 = 1$，$Z_1^{\top}Z_2 = 0$。另外有因变量 $Y$ 的样本值 $Y = (y_1, \ldots, y_n)^{\top}$，也已做好中心化，即 $\sum_{i=1}^n y_i/n = 0$。
\begin{enumerate}
\item (5 pts) 考虑简单回归 $Y \sim Z_1$，由于变量都已中心化，可以不用加入截距。此时最小二乘的目标为
\[
\min_{\gamma_1}\ \|Y - \gamma_1 Z_1\|^2.
\]
证明：此时回归的 $R^2$ 等于 $\rho_1^2$，其中
\[
\rho_1 = \frac{Z_1^{\top}Y}{\|Z_1\|\,\|Y\|}
\]
为 $Z_1$ 与 $Y$ 的样本相关系数。
\item (5 pts) 考虑二元回归 $Y \sim Z_1 + Z_2$，同样不用加入截距。证明：用最小二乘法估计出的两个斜率为
\[
\hat{\gamma}_1 = Z_1^{\top}Y, \qquad \hat{\gamma}_2 = Z_2^{\top}Y.
\]
\item (5 pts) 假设变量 $X_1 = Z_1$，$X_2 = Z_1 + Z_2$，其样本即为 $X_1 = Z_1$，$X_2 = Z_1 + Z_2$。现在用 $X_1, X_2$ 作为自变量，做回归 $Y \sim X_1 + X_2$（不加截距），并用岭回归做正则化，即求解
\[
\min_{\beta_1, \beta_2}\ \|Y - \beta_1 X_1 - \beta_2 X_2\|^2 + \lambda\left(\beta_1^2 + \beta_2^2\right).
\]
请求出岭回归的解，将其写成依赖于 $\hat{\gamma}_1$, $\hat{\gamma}_2$, $\lambda$ 的表达式。
\end{enumerate}

\end{document}
